% =====================================================================
%  IEEE Signal Processing Letters — manuscript
% ---------------------------------------------------------------------
%  Format: IEEEtran `journal` class, 2-column.
%  Length limit: 4 pages typeset content including references, figures,
%  tables, and bibliography. Submissions exceeding 4 pages are returned
%  without review.
%  Style notes per SPL author guidelines:
%   - Compact title; no acronyms in title unless universally understood.
%   - Single contribution; no omnibus.
%   - Abstract: 100-200 words.
%   - 3 to 5 Index Terms (IEEE keyword taxonomy where possible).
%   - References in IEEE numbered style via IEEEtran.bst.
%  Compile:
%      pdflatex main && bibtex main && pdflatex main && pdflatex main
% =====================================================================
\documentclass[journal]{IEEEtran}

\usepackage{cite}
\usepackage{amsmath,amssymb}
\usepackage{graphicx}
\usepackage{booktabs}
\usepackage{microtype}
\usepackage{url}

\hyphenation{Sin-ha-la Ta-mil Vox-Ce-leb}

\newcommand{\dcf}{\mathrm{minDCF}}
\newcommand{\cllr}{C_{\mathrm{llr}}}

\begin{document}

\title{Speaker Verification for Sinhala and Tamil:\\
A Reproducible Benchmark with Zero-Shot and\\
Training-Free Adaptation Baselines}

\author{%
    [First~Author],~\IEEEmembership{Student~Member,~IEEE,}
    [Second~Author],~\IEEEmembership{Member,~IEEE,}
\thanks{Manuscript received [Month] [Day], [Year]; revised [Month] [Day], [Year].}%
\thanks{The authors are with the [Department], [University],
[City], Sri~Lanka (e-mail: [email@institution]).}%
\thanks{Code, configurations, and trial lists are released at
\texttt{github.com/[user]/voxceleb\_trainer-SL-SPV}.}}

\markboth{IEEE Signal Processing Letters,~Vol.~XX, No.~X, [Month]~[Year]}%
{[Author]~\MakeLowercase{\textit{et~al.}}: Speaker Verification for
Sinhala and Tamil}

\maketitle

\begin{abstract}
Sinhala and Tamil, the two official languages of Sri Lanka, have no
published speaker-verification (SV) benchmark. We present the first
reproducible SV benchmark for these languages, built from public
OpenSLR read-speech corpora: 527 speakers (478 Sinhala, 49 Tamil)
with 3{,}000 target and 9{,}000 impostor trials per language,
deterministically generated from a single seed. On this benchmark we
report the first zero-shot baselines for four state-of-the-art
VoxCeleb-trained systems (ECAPA-TDNN and ReDimNet-B1/B2/B6):
equal-error rates of 2.71--4.78\% on Sinhala and 1.48--3.21\% on
Tamil, i.e., $2$--$8\times$ their published VoxCeleb1-O figures even
on clean read speech, with poor minimum detection cost
(0.24--0.48 at $P_{\mathrm{tgt}}{=}0.01$) throughout. A training-free
adaptation study shows that adaptive score normalization with a small
in-domain cohort recovers up to 39\% relative EER (best system:
2.07\%/0.90\%), that PLDA helps or hurts depending on corpus recording
structure, and that swapping calibrators between the two languages
inflates $C_{\mathrm{llr}}$ by up to $5.7\times$, making per-language
calibration mandatory. In-language fine-tuning of WavLM confirms and
completes this picture. Trial lists, code, and configurations are
released.
\end{abstract}

\begin{IEEEkeywords}
Speaker verification, low-resource languages, benchmark,
score normalization, calibration.
\end{IEEEkeywords}

\IEEEpeerreviewmaketitle

% =====================================================================
\section{Introduction}
\label{sec:intro}
% =====================================================================

\IEEEPARstart{S}{peaker} verification (SV) on English benchmarks has
matured: publicly released systems reach 0.4--0.8\% equal-error rate
(EER) on VoxCeleb1-O \cite{desplanques2020ecapa, yakovlev2024redimnet,
nagrani2017voxceleb}. How well these released checkpoints transfer to
unseen low-resource languages is far less well characterized. For the
two official languages of Sri Lanka --- Sinhala (Indo-Aryan,
$\sim$17\,million speakers) and Tamil (Dravidian, $\sim$5\,million Sri
Lankan speakers) --- the gap is complete: neither language appears in
any catalogued SV corpus, and we could find no published SV results
for modern architectures on Sinhala at all.

Evidence from other languages suggests transfer is costly:
VoxCeleb-trained systems degrade by a factor of 2--4 in EER on
CN-Celeb \cite{fan2020cnceleb}, and cross-lingual trials suffer
systematic score shifts that break calibration
\cite{matejka2017analysis, thienpondt2022tackling}. Whether these
findings extend to Sinhala and Tamil, whether the published VoxCeleb
ranking of systems is preserved under the shift, and how much of the
loss is recoverable \emph{without any training} are open questions
with direct deployment relevance.

This letter contributes:
\begin{itemize}
\item \textbf{SL-SPV benchmark v0} --- the first reproducible SV
benchmark for Sinhala and Tamil: 527 speakers drawn from public
OpenSLR corpora \cite{kjartansson2018crowdsourced, he2020opensource},
12{,}000 trials per language constructed deterministically from a
single seed, and evaluation at two operating points
($\dcf$ at $P_{\mathrm{tgt}}=0.01$ and $0.05$) alongside EER.
\item \textbf{Zero-shot baselines} for four state-of-the-art released
checkpoints --- ECAPA-TDNN \cite{desplanques2020ecapa,
ravanelli2021speechbrain} and ReDimNet-B1/B2/B6
\cite{yakovlev2024redimnet} --- quantifying the English$\to$Sinhala/%
Tamil degradation at $2$--$8\times$ published VoxCeleb1-O EER.
\item \textbf{A training-free adaptation study} of adaptive score
normalization (AS-Norm) \cite{matejka2017analysis}, two-covariance
PLDA \cite{sizov2014unifying, kenny2010}, and score calibration,
including the first cross-lingual calibration-transfer measurement
for this language pair. We additionally verify the conclusions
against in-language fine-tuning of WavLM \cite{chen2022wavlm}.
\end{itemize}

Section~\ref{sec:benchmark} describes the benchmark,
Section~\ref{sec:zeroshot} the zero-shot results,
Section~\ref{sec:adapt} the adaptation study,
Section~\ref{sec:limitations} the limitations, and
Section~\ref{sec:conclusion} concludes.

% =====================================================================
\section{The SL-SPV Benchmark (v0)}
\label{sec:benchmark}
% =====================================================================

\subsection{Corpora}
\label{ssec:corpora}

Sinhala audio comes from OpenSLR SLR52
\cite{kjartansson2018crowdsourced}, a crowdsourced read-speech corpus
recorded at 16\,kHz; retaining speakers with at least 10 utterances
yields \textbf{478 speakers} ($\sim$185k utterances). SLR52 publishes
no gender or session metadata. Tamil audio comes from OpenSLR SLR65
\cite{he2020opensource}, a studio-recorded read-speech corpus
(48\,kHz, downsampled to 16\,kHz), yielding \textbf{49 speakers}
(4{,}286 utterances) with speaker gender recoverable from the file
naming convention. Both corpora are CC~BY-SA licensed, so the
benchmark is fully redistributable as trial lists over public audio.

\subsection{Trial protocol}
\label{ssec:protocol}

Each speaker's utterances are split 80/20 into a training/adaptation
pool and a held-out test pool (seed 42); all trials are drawn from
test utterances only. Per language we sample \textbf{3{,}000 target
and 9{,}000 impostor trials} (1:3). Target trials pair same-speaker,
same-language utterances. Impostor trials are same-language and
\emph{same-gender whenever gender is known for both speakers}: the
Tamil list contains zero cross-gender impostor pairs, while Sinhala
impostors are not gender-controlled because SLR52 lacks the metadata
(Section~\ref{sec:limitations}). Pairs are deduplicated as unordered
pairs, self-pairs are excluded, and all sampling is reproducible from
the single seed.

\subsection{Metrics}
\label{ssec:metrics}

We report EER (mean of false-positive and false-negative rates at the
crossing) and $\dcf$ at $P_{\mathrm{tgt}} \in \{0.01, 0.05\}$ with
$C_{\mathrm{miss}}{=}C_{\mathrm{fa}}{=}1$, making the numbers
comparable with both the NIST and VoxCeleb literatures. Calibration
experiments additionally report $\cllr$. The evaluation harness
treats unlabeled trials as a hard error rather than imputing labels.

% =====================================================================
\section{Zero-Shot Baselines}
\label{sec:zeroshot}
% =====================================================================

\begin{table}[!t]
\renewcommand{\arraystretch}{1.12}
\caption{Zero-shot results on SL-SPV v0 (full 12{,}000-trial lists
per language): EER (\%) and $\dcf$ at $P_{\mathrm{tgt}}=0.01/0.05$.
``V1-O'' is each checkpoint's published VoxCeleb1-O EER (\%).}
\label{tab:zeroshot}
\centering
\footnotesize
\setlength{\tabcolsep}{2.6pt}
\begin{tabular}{l r r r c r c}
\toprule
& & & \multicolumn{2}{c}{\textbf{Sinhala}} & \multicolumn{2}{c}{\textbf{Tamil}} \\
\cmidrule(lr){4-5}\cmidrule(lr){6-7}
\textbf{System} & \textbf{Par.} & \textbf{V1-O} & EER & DCF$_{.01/.05}$ & EER & DCF$_{.01/.05}$ \\
\midrule
ECAPA-TDNN \cite{desplanques2020ecapa}  & 20M  & 0.80 & 4.78 & .471/.291 & 3.21 & .482/.207 \\
ReDimNet-B1 \cite{yakovlev2024redimnet} & 2.2M & 0.73 & 3.57 & .306/.209 & 1.67 & .406/.162 \\
ReDimNet-B2 \cite{yakovlev2024redimnet} & 4.7M & 0.52 & 4.17 & .422/.244 & 1.97 & .454/.170 \\
ReDimNet-B6 \cite{yakovlev2024redimnet} & 15M  & 0.37 & \textbf{2.71} & .238/.153 & \textbf{1.48} & .422/.160 \\
\bottomrule
\end{tabular}
\end{table}

We evaluate four released VoxCeleb-trained checkpoints without any
adaptation: the SpeechBrain ECAPA-TDNN \cite{desplanques2020ecapa,
ravanelli2021speechbrain} trained on VoxCeleb1+2
\cite{nagrani2017voxceleb, chung2018voxceleb2}, and the large-margin
fine-tuned ReDimNet-B1/B2/B6 checkpoints \cite{yakovlev2024redimnet}.
Scoring is cosine similarity on $\ell_2$-normalized full-utterance
embeddings, with no score normalization or calibration --- the raw
transfer reference that Section~\ref{sec:adapt} improves on.
Table~\ref{tab:zeroshot} reports the results, the first published SV
baselines for these languages. Four observations:

\emph{(i) Language shift alone costs a multiple, not a margin.}
Relative to each checkpoint's published VoxCeleb1-O EER, the SL-SPV
degradation factor spans $4.9$--$8.0\times$ (Sinhala) and
$2.3$--$4.0\times$ (Tamil): Tamil sits within the $2$--$4\times$ band
that the CN-Celeb literature \cite{fan2020cnceleb} reports for
in-the-wild Chinese, and Sinhala exceeds it, despite v0 being clean,
single-session read speech.

\emph{(ii) The VoxCeleb ranking mostly, but not fully, transfers.}
ReDimNet-B6 is best and ECAPA-TDNN worst on both languages, yet B1
beats B2 on both languages despite B2's better VoxCeleb1-O figure ---
and parameter count does not buy robustness monotonically: B1 (2.2M)
outperforms both B2 (4.7M) and ECAPA-TDNN ($\sim$20M).

\emph{(iii) EER alone overstates how solved Tamil is.} Tamil EER is
uniformly lower (studio audio, 49 speakers), but Tamil
$\dcf_{.01}$ is uniformly \emph{worse} than Sinhala for the ReDimNets
(B6: 0.422 vs.\ 0.238): with few speakers, a handful of hard
impostors dominates the low-false-alarm operating point.

\emph{(iv) Scores are usable but badly scaled.} $\dcf_{.01}$ is poor
everywhere (0.24--0.48) while EERs look deployable --- the classic
signature of uncalibrated scores in a mismatched domain
\cite{matejka2017analysis, thienpondt2022tackling}, and the direct
motivation for the next section.

% =====================================================================
\section{Training-Free Adaptation}
\label{sec:adapt}
% =====================================================================

\subsection{Setup}
\label{ssec:adaptsetup}

All levers in this section require no fine-tuning --- only one
embedding-extraction pass over \emph{training-pool} audio, which is
disjoint from every trial utterance. From the training pool
we build (a) a multi-language cohort of 527 utterances (one per
speaker, both languages, following \cite{matejka2017analysis}) and
(b) a PLDA training set of 10{,}480 utterances ($\le$20 per speaker,
527 speakers). Four backends are compared per model and language:
(1) \emph{cosine} (the Table~\ref{tab:zeroshot} baseline, reproduced
exactly as a sanity check); (2) \emph{AS-Norm}: adaptive symmetric
score normalization using, per trial side, the mean/std of its
top-$K$ ($K{=}300$) cohort similarities,
$s' = \tfrac{1}{2}\!\left[(s-\mu_e)/\sigma_e +
(s-\mu_t)/\sigma_t\right]$;
(3) \emph{PLDA}: two-covariance PLDA \cite{sizov2014unifying,
kenny2010} with LDA to 150 dimensions, centering and length
normalization; and (4) PLDA followed by AS-Norm. We apply the grid to ECAPA-TDNN,
ReDimNet-B1, and ReDimNet-B6 --- the worst, smallest, and best
zero-shot systems.

For calibration, each trial list is split 50/50 (stratified, seed 42)
into calibration and evaluation halves; an affine logistic-regression
mapping (score~$\to$~log-likelihood ratio) is fit per condition.
\emph{Self-cal} fits and evaluates within one language;
\emph{cross-cal} fits on one language and evaluates on the other,
measuring the cross-lingual score shift \cite{thienpondt2022tackling}.

\begin{table}[!t]
\renewcommand{\arraystretch}{1.1}
\caption{Training-free backend adaptation on SL-SPV v0:
EER \% ($\dcf$ at $P_{\mathrm{tgt}}{=}0.01$). Bold: best EER per
system and language.}
\label{tab:backend}
\centering
\small
\setlength{\tabcolsep}{5pt}
\begin{tabular}{l l c c}
\toprule
\textbf{System} & \textbf{Backend} & \textbf{Sinhala} & \textbf{Tamil} \\
\midrule
ECAPA-TDNN  & cosine          & 4.78 (.471) & 3.21 (.482) \\
            & \ \ +AS-Norm    & \textbf{3.19} (.330) & 2.24 (.480) \\
            & PLDA            & 6.00 (.598) & 2.46 (.495) \\
            & \ \ +AS-Norm    & 5.80 (.482) & \textbf{2.06} (.374) \\
\midrule
ReDimNet-B1 & cosine          & 3.57 (.306) & 1.67 (.406) \\
            & \ \ +AS-Norm    & \textbf{2.38} (.245) & \textbf{1.47} (.433) \\
            & PLDA            & 5.34 (.452) & 1.77 (.419) \\
            & \ \ +AS-Norm    & 5.23 (.382) & 1.49 (.297) \\
\midrule
ReDimNet-B6 & cosine          & 2.71 (.238) & 1.48 (.422) \\
            & \ \ +AS-Norm    & \textbf{2.07} (.167) & \textbf{0.90} (.376) \\
            & PLDA            & 3.84 (.293) & 1.11 (.340) \\
            & \ \ +AS-Norm    & 3.56 (.254) & 1.16 (.314) \\
\bottomrule
\end{tabular}
\end{table}

\subsection{Results}
\label{ssec:adaptresults}

\subsubsection{AS-Norm is the single most cost-effective lever}
Table~\ref{tab:backend} shows AS-Norm with the small in-domain cohort
delivers 23--39\% relative EER reduction in five of six
model--language pairs (12\% for B1 on Tamil) --- bracketing the
20--30\% band reported for multilingual score normalization
\cite{matejka2017analysis} --- and improves $\dcf_{.01}$ in five of
six pairs as well (best: B6 Sinhala $0.238\to0.167$; the exception is
B1 Tamil, $0.406\to0.433$). The best training-free
system is \textbf{ReDimNet-B6 + AS-Norm: 2.07\% (Sinhala) / 0.90\%
(Tamil)}, at the cost of one cohort file and no training.

\subsubsection{PLDA is condition-dependent}
In-domain PLDA \emph{hurts} Sinhala EER for every model (B6:
$2.71\to3.84$) yet \emph{helps} Tamil EER for two of three models
(ECAPA: $3.21\to2.46$; B6: $1.48\to1.11$; B1: $1.67\to1.77$), and
PLDA+AS-Norm gives the best Tamil $\dcf_{.01}$ for all three models. Our working hypothesis is a benchmark
artifact: the PLDA fit is dominated by the 478 Sinhala speakers of
near-single-session crowdsourced audio, so its within-speaker
covariance underestimates real channel variability, whereas the
studio-recorded Tamil fits the model better. The claim that PLDA wins
under domain shift \cite{wang2022scoring} should therefore not be
dismissed on this evidence; a multi-session corpus
(Section~\ref{sec:limitations}) is the decisive test.

\subsubsection{Cross-lingual score shift is large and real}
Table~\ref{tab:calib} reports $\cllr$ for ReDimNet-B6 under self- and
cross-language calibration. Swapping calibrators across languages
inflates $\cllr$ by up to $5.7\times$ (Sinhala scored with a
Tamil-fit calibrator under AS-Norm): \emph{per-language calibration
is mandatory} for deployment, extending the cross-lingual score-shift
finding of \cite{thienpondt2022tackling} to this language pair.
Notably, PLDA scores are the most \emph{transportable} across
languages (cross-cal penalty only $1.2$--$2.7\times$) even where
PLDA's EER is worse --- generatively modeled scores appear to live in
a more language-invariant domain than cosine scores.

\begin{table}[!t]
\renewcommand{\arraystretch}{1.1}
\caption{Cross-lingual score shift (ReDimNet-B6): $\cllr$ on the
evaluation half with self-language vs.\ cross-language calibration
(inflation factor in parentheses).}
\label{tab:calib}
\centering
\small
\setlength{\tabcolsep}{3.4pt}
\begin{tabular}{l cc cc}
\toprule
& \multicolumn{2}{c}{\textbf{Sinhala eval}} & \multicolumn{2}{c}{\textbf{Tamil eval}} \\
\cmidrule(lr){2-3}\cmidrule(lr){4-5}
\textbf{Backend} & self & Tamil-fit & self & Sinhala-fit \\
\midrule
cosine        & 0.129 & 0.489 (3.8$\times$) & 0.095 & 0.149 (1.6$\times$) \\
\ \ +AS-Norm  & 0.099 & 0.565 (5.7$\times$) & 0.071 & 0.167 (2.4$\times$) \\
PLDA          & 0.171 & 0.390 (2.3$\times$) & 0.073 & 0.084 (1.2$\times$) \\
\ \ +AS-Norm  & 0.163 & 0.442 (2.7$\times$) & 0.066 & 0.094 (1.4$\times$) \\
\bottomrule
\end{tabular}
\end{table}

\subsection{Verification against in-language fine-tuning}
\label{ssec:finetune}

To check that the training-free conclusions are not an artifact of
frozen encoders, we additionally fine-tuned WavLM-base+
\cite{chen2022wavlm} (attentive-statistics pooling, 192-d embeddings,
AAM-softmax \cite{deng2019arcface}, margin 0.2, scale 30) on
31{,}049 training-pool utterances (all 527 speakers, capped at 60
utterances per speaker). Full fine-tuning reaches
$1.69\pm0.06\%$ (Sinhala) and $1.41\pm0.10\%$ (Tamil) EER over three
seeds --- past the best zero-shot system --- whereas LoRA
\cite{hu2022lora} ($r{=}8$, 1.1\% of parameters) reaches only
$3.73\pm0.49\%$ / $3.11\pm0.61\%$ on this closed-set benchmark.
Composing the backends with the fine-tuned encoder completes the
domain-shift picture: PLDA becomes redundant or harmful on
\emph{both} languages (Sinhala $1.63\to2.46$ at seed 42), exactly as
predicted for in-domain embeddings \cite{wang2022scoring}, while
AS-Norm's EER gain saturates but it still improves the
operating-point metrics ($\dcf_{.01}$ Sinhala $0.180\to0.144$, Tamil
$0.175\to0.149$, with $\cllr$ down $\sim$15\%). Backend adaptation
thus remains worth keeping for deployment even once the encoder is
adapted.

% =====================================================================
\section{Limitations}
\label{sec:limitations}
% =====================================================================

The v0 benchmark has known limitations that any use of its numbers
must respect. (1)~Both corpora are read speech with, in effect, a
single recording session per speaker, so absolute EERs are
optimistic; the numbers are valid for \emph{relative} comparison of
systems and backends, not for comparison against in-the-wild
benchmarks. (2)~The evaluation is closed-set: the adaptation/training
pool contains the same 527 speakers as the trials
(utterance-disjoint), which is standard for domain-adaptation PLDA
but must be stated --- it also removes the overfitting penalty that
usually favors parameter-efficient over full fine-tuning, so the
LoRA result of Section~\ref{ssec:finetune} should not be read as
general. (3)~Sinhala impostor trials are not gender-controlled
(SLR52 has no gender metadata); Tamil trials are. (4)~SLR65 Tamil is
Indian-accented rather than Sri Lankan Tamil. (5)~There is no
telephony/codec condition and no code-switched speech. A v1 rerun on
SLCeleb, a 280-speaker in-the-wild Sri Lankan corpus with
multi-session and bilingual speakers (IEEE DataPort), is planned to
address (1)--(4) and to provide open-set and genuine cross-lingual
trials.

% =====================================================================
\section{Conclusion}
\label{sec:conclusion}
% =====================================================================

This letter establishes the first reproducible speaker-verification
benchmark for Sinhala and Tamil, publishes the first zero-shot
baselines for four state-of-the-art VoxCeleb-trained systems on it,
and characterizes what training-free adaptation can and cannot
repair: AS-Norm with a 527-utterance cohort recovers up to 39\%
relative EER (best system 2.07\%/0.90\%), PLDA's value hinges on
corpus recording structure, and calibration does not transfer across
the two languages. Trial lists, code, and configurations are
released to make every number in this letter reproducible from a
single seed and to serve as the reference point for the planned
in-the-wild v1 benchmark.

% =====================================================================
% References
% =====================================================================
\bibliographystyle{IEEEtran}
\bibliography{../references}

\end{document}
